\documentclass[11pt]{article}
\usepackage{amsmath, amssymb}
\usepackage[margin=1in]{geometry}
\usepackage{booktabs}
\usepackage{siunitx}

\newcommand{\code}[1]{\texttt{#1}}
\DeclareMathOperator*{\argmax}{arg\,max}
\DeclareMathOperator*{\argmin}{arg\,min}

\title{Problem Formulation: Global Design and Daily Operation of the Solar SAWH Device}
\author{}
\date{}

\begin{document}
\maketitle

\section*{Scope}

This document states the optimization problem that \code{sawh\_bayesopt} solves and the
approximations each solver makes. The physics behind the day map $F$ below is in
\code{solar\_lumped/docs/governing\_eq.tex}. Code references are to
\code{sawh\_bayesopt/src/sawh\_bayesopt/}, unless noted otherwise.

The device is the PAM--LiCl hydrogel solar SAWH panel. It runs one cycle per day: it
absorbs at night with the panel open, then it is sealed and desorbs while the sun is up. The
problem is posed once per map cell (one Open-Meteo site), over one weather year (2024,
$D = 366$ days).

\section{Sets, variables and data}

\begin{table}[h]
\centering
\small
\begin{tabular}{@{}llp{0.52\textwidth}@{}}
\toprule
\textbf{Symbol} & \textbf{Domain} & \textbf{Meaning} \\
\midrule
$c \in \mathcal{C}$ & $\approx$ 14.7k cells & Map cell (site), with its own weather year and elevation \\
$d \in \{1,\dots,D\}$ & $D = 366$ & Calendar day \\
$w_{c,d}$ & --- & That day's \SI{15}{\minute} $T$, RH, GHI series (the only channels the physics reads) \\
$x = (H_0, L_g, SL)$ & $\mathcal{X}$ & \textbf{Design}: hydrogel thickness, vapor gap, salt loading \\
$u_d = (\delta_s, \delta_o, \theta)$ & $\mathcal{U}$ & \textbf{Control}: seal offset, open offset, tilt \\
$s_d = (c_w, H)$ & $\mathbb{R}_{>0}^2$ & \textbf{State}: gel water concentration and thickness at the start of day $d$ \\
$\phi_{c,d}$ & $\mathbb{R}^{k+7}$ & Day features (Section~\ref{sec:surrogate}) \\
\bottomrule
\end{tabular}
\end{table}

\paragraph{Design box} (\code{design\_space.DesignBounds}, read from the \code{parameters.xlsx} sweep columns):
\begin{equation}
\mathcal{X} = [\SI{1}{\milli\metre}, \SI{10}{\milli\metre}] \times
[\SI{7}{\milli\metre}, \SI{60}{\milli\metre}] \times [1, 8].
\end{equation}
The \SI{7}{\milli\metre} gap floor is a transport wall, not a soft preference. The box's
worst corner is \SI{6.03}{\milli\metre} of dry gel against that \SI{7}{\milli\metre} minimum
gap, so no design starts with the gel touching the condenser.

\paragraph{Control grid} (\code{daily\_surrogate}):
\begin{equation}
\mathcal{U} = \mathcal{S} \times \mathcal{O} \times \Theta, \qquad
\mathcal{S} = \mathcal{O} = \{-4, -3.75, \dots, 4\}\,\si{\hour}, \qquad
\Theta = \{0, 5, \dots, 60\}\si{\degree},
\end{equation}
which gives $|\mathcal{S}\times\mathcal{O}| = 33^2 = 1089$ schedules per tilt level and
$|\mathcal{U}| = 14{,}157$. The offsets move the seal and open times relative to GHI sunrise
and sunset, where GHI crosses the night threshold. Negative $\delta_s$ seals before dawn;
positive $\delta_o$ stays shut past dusk. Both edges are clamped so that each phase keeps
at least one sample (\code{solar\_lumped.weather.\_shifted\_desorption\_mask}). The
\SI{0.25}{\hour} step equals the weather time step, so a finer offset would not change the
physics. Tilt enters twice: through plane-of-array transposition of GHI, and through
Hollands $\cos\theta$ in the vapor gap.

\paragraph{Held fixed} (not decision variables): LiCl salt; Case 2 optics and IR
emissivities; insulation gap \SI{5}{\milli\metre}; condenser fin area ratio $A_r = 7.1$
(\code{design\_space.SIMPLE\_FIXED}); $h_{amb}$; all financial parameters
(\code{LCOEconomicParams}). Each cell's real ambient pressure comes from its elevation.

\section{The day map and the year}

One simulated day is a map, from the JAX physics (\code{solar\_lumped/gpu\_sweep/jax\_daily\_cycle}):
\begin{equation}
\bigl(m_d,\; s_{d+1},\; \kappa_d\bigr) = F\bigl(x, u_d, s_d;\, w_{c,d}\bigr),
\label{eq:daymap}
\end{equation}
where $m_d$ is the water collected that day (\si{\kilogram\per\metre\squared}) and
$\kappa_d \in \{0,1\}$ flags a gel that reached the swelling ceiling (the gel--condenser
clearance) during the day. The gel carries its state from one day to the next, so the year
is a chained trajectory. Its start is the cyclic steady state of day 1:
\begin{equation}
s_1 = s^\ast \quad\text{with}\quad s^\ast = F_s\bigl(x, u_1, s^\ast;\, w_{c,1}\bigr),
\label{eq:cyclic}
\end{equation}
which the physics finds by Aitken $\Delta^2$ extrapolation. The annual figures of merit are
\begin{equation}
\bar m(x, u_{1:D}) = \frac{1}{D}\sum_{d=1}^{D} m_d,
\qquad
\bar\kappa(x, u_{1:D}) = \frac{1}{D}\sum_{d=1}^{D} \kappa_d .
\end{equation}

\section{Objective: levelized cost of water}

The cost per m$^2$ of panel is (\code{solar\_lumped.economics.lcow\_from\_daily\_yield}):
\begin{equation}
\mathrm{LCOW}(x, \bar m) =
\frac{\bigl(\mathrm{CRF}\cdot f_{TI} + f_{OM}\bigr)\,C_{sys} + C_{sorb}(x)}
{f_{util}\cdot 365\,\bar m/\rho_w}
\quad [\text{USD/m}^3],
\label{eq:lcow}
\end{equation}
\begin{equation}
\mathrm{CRF} = \frac{i(1+i)^L}{(1+i)^L - 1},
\qquad
C_{sorb}(x) = \left[\frac{p_{salt}\,SL + p_{poly}}{1 + SL} + r_w\,p_{w}\right]
\frac{\rho_{dry}\,H_0}{L_{gel}} .
\label{eq:csorb}
\end{equation}
The symbols are: $C_{sys}$, the flat system bill of materials; $f_{TI}$, the
total-investment factor; $f_{OM}$, the maintenance fraction; $f_{util}$, utilization; $i$
and $L$, the discount rate and system life; and $L_{gel}$, the gel replacement interval.
$C_{sorb}$ is the annualized cost of re-casting the gel. It is the dry composite mass
$\rho_{dry}H_0$, priced at the salt/polymer blend plus the DI water in the recipe.

The design enters LCOW in two places, which is what creates the trade-offs:
\begin{itemize}
\item \textbf{Through the yield} $\bar m$: thicker and saltier gels hold more water, but
they desorb more slowly and swell further.
\item \textbf{Through $C_{sorb}$}: cost is linear in $H_0$, and grows with $SL$ when
$p_{salt} > p_{poly}$.
\end{itemize}
The vapor gap and the controls act only through $\bar m$, since neither is priced.

\textbf{For fixed $x$, LCOW is strictly decreasing in $\bar m$.} This property is what
justifies the decomposition in Section~\ref{sec:decomp}.

\section{The full problem}

For each cell $c$, choose a design and a control policy:
\begin{equation}
\begin{aligned}
\min_{x \in \mathcal{X},\; \pi \in \Pi} \quad & \mathrm{LCOW}\bigl(x, \bar m(x, u_{1:D})\bigr) \\
\text{s.t.} \quad & (m_d, s_{d+1}, \kappa_d) = F(x, u_d, s_d; w_{c,d}), \quad d = 1,\dots,D, \\
& s_1 \text{ cyclic (Eq.~\eqref{eq:cyclic})}, \\
& u_d = \pi_d(\mathcal{I}_d) \in \mathcal{U}, \\
& \bar\kappa(x, u_{1:D}) \le \kappa_{\max} = 0.05 .
\end{aligned}
\label{eq:full}
\end{equation}

\paragraph{Swelling constraint.} A design is infeasible if its gel sits on the swelling
ceiling on more than 5\% of days (\code{two\_stage.CAPPED\_DAY\_FRACTION\_MAX}). On those
days, uptake is set by the clearance rather than by the isotherm; in the real device, this
is brine reaching the condenser. The constraint depends on the climate, and it binds mainly
at high RH combined with high $SL$. A share-of-days threshold is a blunt stand-in for a
leakage model, which the repo does not have. Designs with no finite positive yield are also
infeasible. Infeasible designs are scored at
$\mathrm{LCOW}_{pen} = 10^4$~USD/m$^3$.

\paragraph{Policy class $\Pi$.} The policy class has two independent axes. The \emph{tilt
mode} says how often $\theta$ may change:
\begin{align}
\text{fixed:}    &\quad \theta_d = \theta \quad \forall d \quad (\theta \text{ is effectively a design variable}), \\
\text{seasonal:} &\quad \theta_d = \theta_{q(d)}, \quad q(d) \in \{1,2,3,4\} \text{ the calendar quarter}, \\
\text{daily:}    &\quad \theta_d \text{ free}.
\end{align}
These modes nest, so the LCOW gain at each step up is the value of moving the panel more
often.

The \emph{information regime} sets what the controller knows, $\mathcal{I}_d$, when it picks
$u_d$. It sees decision features $\tilde\phi_{c,d}$ in place of the true $\phi_{c,d}$:
\begin{align}
\text{hindsight:}      &\quad \tilde\phi_{c,d} = \phi_{c,d} \quad \text{(perfect foresight; an upper bound)}, \\
\text{persistence:}    &\quad \tilde\phi_{c,d} = \phi_{c,d-1} \quad \text{(day-ahead forecast stand-in; day 1 uses day } D), \\
\text{climatological:} &\quad \tilde\phi_{c,d} = \tfrac{1}{|M(d)|}\textstyle\sum_{d' \in M(d)} \phi_{c,d'} \quad \text{(mean day of its month)}, \\
\text{constant:}       &\quad u_d = u \quad \forall d \quad \text{(one annual schedule; fixed tilt only)}.
\end{align}
A control is always \emph{chosen} on $\tilde\phi$ and \emph{scored} on the real day
$\phi_{c,d}$. The constant regime is the control space of the per-site reference BO
(Section~\ref{sec:reference}).

\section{Decomposition: design per cell, control per day}
\label{sec:decomp}

Since LCOW decreases in $\bar m$ for fixed $x$, the inner problem can be posed in water
rather than dollars:
\begin{equation}
\min_{x \in \mathcal{X}} \; \mathrm{LCOW}\bigl(x, \bar m^\ast(x)\bigr)
\;\;\text{s.t.}\;\; \bar\kappa\bigl(x, u^\ast_{1:D}(x)\bigr) \le \kappa_{\max},
\qquad
u^\ast_{1:D}(x) \in \argmax_{\pi \in \Pi} \bar m(x, u_{1:D}),
\label{eq:bilevel}
\end{equation}
with $\bar m^\ast(x) = \bar m(x, u^\ast_{1:D}(x))$. The swelling constraint couples the two
levels, but the inner level ignores it. Controls maximize water, and feasibility is checked
only at the design level. So Eq.~\eqref{eq:bilevel} is a restriction of
Eq.~\eqref{eq:full}: a design that would become feasible under a slightly lower-yield,
less-swelling schedule is rejected rather than rescued.

Keeping the inner level in physical units has a practical benefit. Every map can be
regenerated under new cost assumptions by re-applying Eq.~\eqref{eq:lcow}, with no
retraining.

\section{Surrogate day map}
\label{sec:surrogate}

The inner level evaluates thousands of controls per day, per design, per cell, which rules
out calling $F$ directly. It uses a learned day map instead (\code{daily\_surrogate}):
\begin{equation}
g_\psi(x, u, s, \phi) \;\to\; \bigl(\hat m,\; \hat s',\; \hat p_\kappa\bigr) \approx
\bigl(m_d,\; s_{d+1},\; \Pr[\kappa_d = 1]\bigr).
\end{equation}
$g_\psi$ is an ensemble of $M = 5$ multilayer perceptrons. They share a trunk and have three
heads: $\log(1+m)$, the end-of-day state, and a swelling-cap logit. Each member is trained
on a Poisson(1) bootstrap over \emph{cells}, so the ensemble spread $\hat\sigma_m$ reads as
``unsure about climates like this one''. The mean across members is used everywhere. The
predicted state is clipped to the range seen in training, and
$\hat\kappa_d = \mathbb{1}[\hat p_\kappa > 0.5]$.

\paragraph{Day features.} $\phi_{c,d}$ joins two blocks:
\begin{itemize}
\item the leading $k$ principal components of the day's standardized hourly
$[T, \mathrm{RH}, \mathrm{GHI}]$ vector ($3\times24 = 72$ numbers; the PCA is fitted over
cached days worldwide);
\item seven scalars: RH at the diurnal $T$ minimum, the daytime GHI integral, the
dewpoint depression at peak sun, the diurnal $T$ amplitude, elevation, $|\text{latitude}|$,
and noon zenith.
\end{itemize}
The last two scalars carry the sun path, which tilt acts through. The \code{validate-features}
stage sets $k$: it runs the physics on real years and on years rebuilt from $k$ components,
and compares the yields.

\paragraph{Training data.} Representative cells are chosen by KMeans in climate-descriptor
space (800 clusters), with quotas topping up the high-altitude, hyper-arid, monsoonal and
coastal-humid regimes. Each cell gets 16 scrambled-Sobol designs. Each (cell, design) is
simulated as one chained physics year, with $u_d$ drawn uniformly from $\mathcal{U}$
\emph{independently every day}. One year therefore yields $\sim$365 rows that cover the
control space, and the states in those rows are the ones real operation produces. Each
year's first day is dropped, because its start state comes from the Aitken extrapolation
rather than from a simulated previous day. About 15\% of cells per regime are held out
\emph{whole}: a split by row would leak each climate across train and test.

\section{Solution method}

\subsection{Inner level: enumeration on the surrogate}

There is no gradient step and no local-minimum risk, and the result is a discrete schedule,
which is what a real controller runs (\code{two\_stage}). Let $\mathcal{U}_\theta$ denote
the 1089 schedules at tilt $\theta$. The \textbf{greedy daily walk} is
\begin{align}
u_d &= \argmax_{u \in \mathcal{U}_d} \hat m\bigl(x, u, \hat s_d, \tilde\phi_{c,d}\bigr), \\
(\hat m_d, \hat s_{d+1}) &= g_\psi\bigl(x, u_d, \hat s_d, \phi_{c,d}\bigr),
\end{align}
where the candidate set $\mathcal{U}_d$ depends on the mode:
\begin{itemize}
\item \textbf{daily tilt:} $\mathcal{U}_d = \mathcal{U}$;
\item \textbf{fixed tilt:} one greedy walk per $\theta \in \Theta$ over $\mathcal{U}_\theta$,
keeping the $\theta$ with the largest $\sum_d \hat m_d$;
\item \textbf{seasonal tilt:} the fixed-tilt procedure applied quarter by quarter, with
each quarter's end state chained into the next.
\end{itemize}
The \textbf{constant} regime is solved exactly. Every $u \in \mathcal{U}$ is chained through
the whole year on $g_\psi$, and the one with the largest $\sum_d \hat m_d$ is kept. The
cyclic start (Eq.~\eqref{eq:cyclic}) is approximated by 8 rounds of the day-1 greedy step
from the physics' initial loading.

The greedy walk is \textbf{myopic}: it ignores what today's choice does to tomorrow's
starting state. Dynamic programming over a discretized state would be the exact version of
Eq.~\eqref{eq:bilevel}. The \code{validate-bo} gap measures what the myopia costs.

\subsection{Outer level: constrained Bayesian optimization over $\mathcal{X}$}

For each cell, the outer level minimizes $J_c(x) = \mathrm{LCOW}(x, \hat{\bar m}^\ast(x))$,
or $\mathrm{LCOW}_{pen}$ if $x$ is infeasible. The setup:
\begin{itemize}
\item \textbf{Models:} a GP regressor with a Mat\'ern-5/2 plus white-noise kernel, fitted on
the unit cube to feasible points only, and a GP classifier for $P(\text{feasible}\mid x)$.
\item \textbf{Acquisition:} constrained expected improvement, maximized by differential
evolution:
\begin{equation}
\alpha(x) = \mathrm{EI}(x)\cdot P(\text{feasible}\mid x),
\qquad
\mathrm{EI}(x) = (y^\ast - \mu - \xi)\,\Phi(z) + \sigma\,\varphi(z),
\quad z = \frac{y^\ast - \mu - \xi}{\sigma}.
\end{equation}
\item \textbf{Batches:} 4 proposals per round, via Kriging-Believer.
\item \textbf{Budget:} 8 Sobol initial points and 24 evaluations in total for the training
(``anchor'') cells. Every other cell is warm-started from the optima of its 3 nearest
solved anchors in standardized climate-descriptor space, with a budget of 12.
\end{itemize}
Cells advance in lockstep, so each round costs one vmapped surrogate sweep across all cells.

\subsection{Active refinement}

The surrogate needs to be accurate near each cell's optimum, not uniformly over the box. So
after a first solve, the pipeline runs further physics on selected cells:
\begin{itemize}
\item cells are drawn with probability $\propto \hat\sigma_m/\hat m$ at their optimum;
\item each drawn cell's optimum and 4 perturbations (within $\pm10\%$ of each span) are run
under the surrogate's own schedule, with jitter;
\item $g_\psi$ is then refit. This takes one or two rounds.
\end{itemize}

\section{Reference problem: per-site BO on true physics}
\label{sec:reference}

The validation reference (\code{bayesopt}, \code{evaluator}) solves Eq.~\eqref{eq:full}
directly, with no surrogate, under the constant regime with fixed tilt. Tilt and the
schedule move into the design vector:
\begin{equation}
\min_{x' \in \mathcal{X}'} \mathrm{LCOW}\bigl(x', \bar m_F(x')\bigr),
\qquad
x' = (H_0, L_g, SL, \theta, \delta_s, \delta_o),
\end{equation}
with $\theta \in [0, 60]\si{\degree}$ continuous, offsets snapped to the \SI{0.25}{\hour}
grid, and $\bar m_F$ the true-physics annual mean. It uses the same GP, constrained EI and
Kriging-Believer code: 24 Latin-hypercube initial points, a budget of 50, batches of 3, and
it stops after 3 rounds with less than 0.5\% improvement.

The \code{validate-bo} stage scores the two-stage design and schedule with true physics on
held-out cells. The headline number is the gap
\begin{equation}
\Delta_c = \mathrm{LCOW}_F\bigl(x_c^{2s}, u^{2s}_{c,1:D}\bigr) - \mathrm{LCOW}_F\bigl(x_c'^{\,BO}\bigr)
\quad [\text{USD/m}^3].
\end{equation}
Under (fixed, constant), $\Delta_c$ is a like-for-like optimality gap. Under the daily
modes, it also includes the value of daily control.

\paragraph{Differences between the two problems that $\Delta_c$ also absorbs:}
\begin{itemize}
\item \textbf{Tilt:} continuous in the reference, but on a \SI{5}{\degree} grid in the
pipeline.
\item \textbf{Swelling constraint:} the reference only \emph{records} the swelling-cap
flag, and its feasibility is a finite positive yield and LCOW. The
$\bar\kappa \le 0.05$ constraint is enforced only by the pipeline. So the reference
optimizes over a feasible set that is at least as large, which biases $\Delta_c$ upward
for cells where the constraint binds.
\end{itemize}

\section*{Summary of approximations}

\begin{table}[h]
\centering
\small
\begin{tabular}{@{}p{0.30\textwidth}p{0.62\textwidth}@{}}
\toprule
\textbf{Approximation} & \textbf{Where it is measured or bounded} \\
\midrule
$F \approx g_\psi$ & \code{holdout\_report.json}: daily and chained-annual error on held-out cells, by regime; design ranking and response signs \\
Weather $\approx$ $k$ PCs + 7 scalars & \code{featval\_summary.csv}: physics on real vs.\ reconstructed years \\
Inner level greedy, not DP & \code{validate-bo} gap under daily modes \\
Inner level ignores the swelling constraint & Not measured; restriction of Eq.~\eqref{eq:full} \\
Cyclic start by 8 surrogate rounds & Not separately measured \\
Climates outside training & Flagged: Mahalanobis distance above the training cells' 99th percentile (\code{ood} column in the maps) \\
\bottomrule
\end{tabular}
\end{table}

\end{document}
